% TOP_PROBLEM: {"id": 2837, "title": "Kirby Problem 3.39", "category_id": 7, "category": {"id": 7, "name": "topology", "display_name": "Topology", "description": "Properties preserved under continuous deformations.", "slug": "topology", "order_index": 7, "created_at": "2026-07-31T15:26:25.671Z"}, "status": "open", "problem_number": "KP-3.39", "set_id": 11, "statusQualification": "Parts (a) and (b) are resolved in Chen–Lodha arXiv:2510.26073v2; this definition retains the open unboundedness targets (c) and (d)."}
% FIRST_POSTED: 2026-10-01
% ENTRY_KIND: statement_only
% UnsolvedMath ID 2837; source catalogue code KP-3.39
% !TeX program = lualatex
% Definition and sources only; this file is not an LLM solution attempt.
\documentclass[11pt,a4paper]{article}
\usepackage[margin=25mm]{geometry}
\usepackage{fontspec}
\setmainfont{lmroman10-regular.otf}[BoldFont=lmroman10-bold.otf,
 ItalicFont=lmroman10-italic.otf,BoldItalicFont=lmroman10-bolditalic.otf]
\usepackage{amsmath,amssymb,mathtools}
\usepackage{unicode-math}
\setmathfont{latinmodern-math.otf}
\AtBeginDocument{\let\setminus\smallsetminus}
\usepackage{xurl,hyperref}
\hypersetup{colorlinks=true,urlcolor=blue}
\setcounter{secnumdepth}{0}
\setlength{\parindent}{0pt}
\setlength{\parskip}{5pt}
\setlength{\emergencystretch}{3em}
\begin{document}
\section*{Kirby Problem 3.39}\label{2837}
{\small UnsolvedMath ID 2837 · KP-3.39 · Topology · Kirby's Problems in Low-Dimensional Topology}
\subsection{Definitions and mathematical statement}
The weight $w(G)$ of a finitely generated group is the smallest size of a
set $A\subseteq G$ whose normal closure is $G$:
\[
 G=\langle\!\langle A\rangle\!\rangle
   =\langle gag^{-1}:g\in G,\ a\in A\rangle.
\]
Set $w(1)=0$. A group is perfect if $G=[G,G]$, or equivalently its
abelianization is trivial. An integral homology three-sphere is a closed
connected oriented manifold $Y$ with $H_*(Y;\mathbb Z)\cong H_*(S^3;\mathbb Z)$;
its fundamental group is perfect.

The catalogue asks four questions: (a) does every nontrivial finitely generated
perfect group have weight one? (b) can an integral homology three-sphere
group have weight greater than one? (c) for every $N$, can such a group have
weight greater than $N$? (d) are the weights unbounded already for Seifert
fibered homology spheres $\Sigma(a_1,\ldots,a_r)$? These are circle-fibered
three-manifolds over a sphere with exceptional fiber multiplicities $a_i$,
pairwise coprime in the homology-sphere setting.

\textbf{Status distinction.} Chen--Lodha's 2025 preprint gives a negative
answer to (a) and constructs hyperbolic integral homology spheres answering
(b) positively [S3]. The remaining targets recorded here are the unboundedness
questions (c) and (d); the lower bound two from that result does not establish
arbitrarily large weight.
\subsection{Short English statement}
Determine how large the normal generating number of an integral homology three-sphere group can be, particularly among Seifert fibered examples.
\subsection{Sources}
\begin{itemize}
\item[S1] ULAM.AI and the UnsolvedMath Contributors, supplied problems.json, record 2837 (KP-3.39), Kirby's Problems in Low-Dimensional Topology. Catalogue identity and source transcription; CC BY 4.0.
\url{https://huggingface.co/datasets/ulamai/UnsolvedMath}.
\item[S2] K3: A New Problem List in Low-Dimensional Topology, Problem 3.39. Expanded formulation of the supplied catalogue statement.
\url{https://math.berkeley.edu/sites/default/files/surv-295-ruberman-watermarked-author-pdf.pdf}.
\item[S3] L. Chen and Y. Lodha, The Wiegold problem and free products of left-orderable groups, v2 (2025), Theorems A and C. Resolves parts (a) and (b), distinguished from the remaining questions.
\url{https://arxiv.org/abs/2510.26073v2}.
\end{itemize}
\end{document}
